\documentclass[11pt,a4paper]{article}
\usepackage{fontspec,amsmath,amssymb,amsthm,microtype}
\setmainfont{lmroman10-regular.otf}[BoldFont=lmroman10-bold.otf,ItalicFont=lmroman10-italic.otf,BoldItalicFont=lmroman10-bolditalic.otf]
\usepackage{xcolor}
\usepackage[most]{tcolorbox}
\usepackage{mathrsfs}
\providecommand{\lt}{<}
\providecommand{\gt}{>}
\setmonofont{DejaVuSansMono.ttf}[BoldFont=DejaVuSansMono-Bold.ttf,ItalicFont=DejaVuSansMono-Oblique.ttf,BoldItalicFont=DejaVuSansMono-BoldOblique.ttf,Scale=0.86]
\newfontfamily\leanfallback{STIXTwoMath-Regular.otf}[Scale=0.86]
\definecolor{leanrule}{gray}{0.62}
\definecolor{leanback}{gray}{0.962}
\definecolor{leantitle}{gray}{0.22}
\definecolor{leanmuted}{gray}{0.38}
\newtcolorbox{leanbox}[1]{enhanced,breakable,colback=leanback,colframe=leanrule,
  boxrule=0.4pt,arc=0.6mm,left=2.4mm,right=2.4mm,top=1.8mm,bottom=1.8mm,
  before skip=0.8em,after skip=1em,fonttitle=\small\bfseries,coltitle=leantitle,
  colbacktitle=leanback,titlerule=0.3pt,title={#1}}
\newlength{\leancharwidth}
\newenvironment{leancode}{\par\smallskip\begingroup\ttfamily\footnotesize\raggedright
  \settowidth{\leancharwidth}{M}\setlength{\parskip}{0pt}\setlength{\parindent}{0pt}
  \frenchspacing}{\par\endgroup\smallskip}
\newcommand{\leanline}[2]{\hangindent=\dimexpr#1\leancharwidth+4\leancharwidth\relax
  \hangafter=1\noindent\hspace*{#1\leancharwidth}#2\par}
\newcommand{\leansource}[1]{\par{\footnotesize\color{leanmuted}#1}\par}
\newcommand{\leanlabel}[1]{\par\smallskip\noindent{\small\bfseries #1}\par}
\usepackage[margin=25mm]{geometry}
\usepackage[unicode,hidelinks]{hyperref}
\hypersetup{pdfauthor={LangevinLean Project}}
\setlength{\emergencystretch}{3em}
\setlength{\parskip}{0.4em}
\setlength{\parindent}{0pt}
\allowdisplaybreaks[2]
\raggedbottom
\title{\texorpdfstring{Equality of weighted \(L^{2}\) norms and Euclidean phase coupling costs}{Equality of weighted L² norms and Euclidean phase coupling costs}}
\author{LangevinLean Project}
\date{}
\begin{document}\maketitle
Draft mathematical exposition; these notes may contain mistakes.


A local estimate in a pair of \(L^{2}\) spaces must refer to the same Euclidean phase geometry as the vector sampler. This section constructs the explicit \(L^{2}\) Euclidean pair and proves its exact norm and representative identities. It then proves that the weighted representative cost is precisely the weighted phase norm used by the local analysis.



The weighted Hilbert-space norm is proved equal to the second-mean finite-dimensional phase cost. This equality converts analytic local-error bounds directly into transport costs of coupled outputs.


\section{\texorpdfstring{Construct a Euclidean phase variable in \(L^{2}\)}{Construct a Euclidean phase variable in L²}}


The coordinate injections \(u\mapsto(u,0)\), \(w\mapsto(0,w)\) are continuous linear maps into the Euclidean product Hilbert space. Lift each through the bounded pullback map and add their outputs. Almost-everywhere identities for lifted linear maps and \(L^{2}\) addition show that the resulting variable represents exactly \((A(\omega),B(\omega))\). Expand its \(L^{2}\) inner product as an integral, use the Euclidean product inner product pointwise, and split the two integrable self-inner-products. This proves \(\|\operatorname{pair}(A,B)\|_{L^2}^2=\|A\|_{L^2}^2+\|B\|_{L^2}^2\). Thus the construction uses the Euclidean geometry, not the default maximum norm of an ordinary product.



Combine velocity and position as one Euclidean phase-valued random variable. Its second-mean squared norm is the sum of the two component second moments. This identifies the functional-analytic phase cost with the finite-dimensional cost used in transport.

\begin{leanbox}{Lean: \texttt{phaseInl}, \texttt{phaseInr}, \texttt{lpEuclideanPair} and 2 more}
\leansource{\texttt{phaseInl} \quad KineticLangevin/\allowbreak{}LpWeightedGeometry.lean, lines 20–21}
\begin{leancode}
\leanline{0}{def\ phaseInl\ :\ E\ →L[ℝ]\ WithLp\ 2\ (E\ ×\ E)\ :=}
\leanline{2}{(WithLp.prodContinuousLinearEquiv\ 2\ ℝ\ E\ E).symm.toContinuousLinearMap.comp\ (ContinuousLinearMap.inl\ ℝ\ E\ E)}
\end{leancode}
\leansource{\texttt{phaseInr} \quad KineticLangevin/\allowbreak{}LpWeightedGeometry.lean, lines 23–24}
\begin{leancode}
\leanline{0}{def\ phaseInr\ :\ E\ →L[ℝ]\ WithLp\ 2\ (E\ ×\ E)\ :=}
\leanline{2}{(WithLp.prodContinuousLinearEquiv\ 2\ ℝ\ E\ E).symm.toContinuousLinearMap.comp\ (ContinuousLinearMap.inr\ ℝ\ E\ E)}
\end{leancode}
\leansource{\texttt{lpEuclideanPair} \quad KineticLangevin/\allowbreak{}LpWeightedGeometry.lean, lines 26–27}
\begin{leancode}
\leanline{0}{def\ lpEuclideanPair\ (A\ B\ :\ Lp\ E\ 2\ P)\ :\ Lp\ (WithLp\ 2\ (E\ ×\ E))\ 2\ P\ :=}
\leanline{2}{(phaseInl\ (E\ :=\ E)).compLpL\ 2\ P\ A\ +\ (phaseInr\ (E\ :=\ E)).compLpL\ 2\ P\ B}
\end{leancode}
\leansource{\texttt{lpEuclideanPair\_\allowbreak{}coe} \quad KineticLangevin/\allowbreak{}LpWeightedGeometry.lean, lines 29–39}
\begin{leancode}
\leanline{0}{theorem\ lpEuclideanPair\_coe\ (A\ B\ :\ Lp\ E\ 2\ P)\ :}
\leanline{4}{lpEuclideanPair\ P\ A\ B\ =ᵐ[P]\ fun\ ω\ ↦\ WithLp.toLp\ 2\ (A\ ω,B\ ω)}
\end{leancode}
\leansource{\texttt{lpEuclideanPair\_\allowbreak{}norm\_\allowbreak{}sq} \quad KineticLangevin/\allowbreak{}LpWeightedGeometry.lean, lines 41–52}
\begin{leancode}
\leanline{0}{theorem\ lpEuclideanPair\_norm\_sq\ (A\ B\ :\ Lp\ E\ 2\ P)\ :}
\leanline{4}{‖lpEuclideanPair\ P\ A\ B‖\ \textasciicircum{}\ 2\ =\ ‖A‖\ \textasciicircum{}\ 2\ +\ ‖B‖\ \textasciicircum{}\ 2}
\end{leancode}
\leanlabel{How the Lean proof follows the argument}
\texttt{phase\allowbreak{}Inl} names the mathematical object defined above. Its defining expression appears in the adjacent statement; subsequent identities may unfold this definition directly.




\texttt{phase\allowbreak{}Inr} names the mathematical object defined above. Its defining expression appears in the adjacent statement; subsequent identities may unfold this definition directly.




\texttt{lp\allowbreak{}Euclidean\allowbreak{}Pair} names the mathematical object defined above. Its defining expression appears in the adjacent statement; subsequent identities may unfold this definition directly.




\texttt{lp\allowbreak{}Euclidean\allowbreak{}Pair\_\allowbreak{}coe} gives the formal counterpart of this step. Named proof applications include \texttt{With\allowbreak{}Lp.\allowbreak{}of\allowbreak{}Lp\_\allowbreak{}injective}; the adjacent proof shows their hypotheses and order of use.




\texttt{lp\allowbreak{}Euclidean\allowbreak{}Pair\_\allowbreak{}norm\_\allowbreak{}sq} gives the formal counterpart of this step. Named proof applications include \texttt{integral\_\allowbreak{}congr\_\allowbreak{}ae}; the adjacent proof shows their hypotheses and order of use.
\end{leanbox}

\section{\texorpdfstring{Lift the weighted transform without changing its norm}{Lift the weighted transform without changing its norm}}


Apply the Euclidean pair construction to \(X+bV\) and \(\sqrt{a-b^2}V\). The representative is almost everywhere the original vector weighted transform \(T(V(\omega),X(\omega))\), with all addition and scalar-multiplication representatives justified explicitly. The squared norm identity for the pair is exactly the defining squared norm of the weighted phase transform on the \(L^{2}\) Hilbert space. Both norms are nonnegative, so equality of their squares gives equality of the norms. Finally the standard \(L^{2}\) norm definition and finiteness of an \(L^{2}\) variable give

\[\|T(V(\cdot),X(\cdot))\|_{L^2}=\|(V,X)\|_E.\]

The identity retains extended-real \(L^{2}\) norms on the left and converts the finite norm with the nonnegative extended-real embedding. Consequently the local Hilbert-space norm equals the transport cost of this particular coupling.



The weighted transformation is bounded linear, so applying it before or after passage to square-integrable classes gives the same element. Its norm identity therefore transfers exactly to random phase errors. No extra coordinatewise inequality is needed.

\begin{leanbox}{Lean: \texttt{lpWeightedTransform}, \texttt{lpWeightedTransform\_\allowbreak{}coe}, \texttt{lpWeightedTransform\_\allowbreak{}norm} and 1 more}
\leansource{\texttt{lpWeightedTransform} \quad KineticLangevin/\allowbreak{}LpWeightedGeometry.lean, lines 54–55}
\begin{leancode}
\leanline{0}{def\ lpWeightedTransform\ (a\ b\ :\ ℝ)\ (S\ :\ Lp\ E\ 2\ P\ ×\ Lp\ E\ 2\ P)\ :\ Lp\ (WithLp\ 2\ (E\ ×\ E))\ 2\ P\ :=}
\leanline{2}{lpEuclideanPair\ P\ (S.2\ +\ b\ •\ S.1)\ (Real.sqrt\ (a\ -\ b\ \textasciicircum{}\ 2)\ •\ S.1)}
\end{leancode}
\leansource{\texttt{lpWeightedTransform\_\allowbreak{}coe} \quad KineticLangevin/\allowbreak{}LpWeightedGeometry.lean, lines 57–67}
\begin{leancode}
\leanline{0}{theorem\ lpWeightedTransform\_coe\ (a\ b\ :\ ℝ)\ (S\ :\ Lp\ E\ 2\ P\ ×\ Lp\ E\ 2\ P)\ :}
\leanline{4}{lpWeightedTransform\ P\ a\ b\ S\ =ᵐ[P]\ fun\ ω\ ↦\ weightedPhaseTransform\ a\ b\ (lpPhaseCoe\ P\ S\ ω)}
\end{leancode}
\leansource{\texttt{lpWeightedTransform\_\allowbreak{}norm} \quad KineticLangevin/\allowbreak{}LpWeightedGeometry.lean, lines 69–76}
\begin{leancode}
\leanline{0}{theorem\ lpWeightedTransform\_norm\ (a\ b\ :\ ℝ)\ (S\ :\ Lp\ E\ 2\ P\ ×\ Lp\ E\ 2\ P)\ :}
\leanline{4}{‖lpWeightedTransform\ P\ a\ b\ S‖\ =\ weightedPhaseNorm\ a\ b\ S}
\end{leancode}
\leansource{\texttt{eLpNorm\_\allowbreak{}weightedPhaseCoe} \quad KineticLangevin/\allowbreak{}LpWeightedGeometry.lean, lines 78–83}
\begin{leancode}
\leanline{0}{theorem\ eLpNorm\_weightedPhaseCoe\ (a\ b\ :\ ℝ)\ (S\ :\ Lp\ E\ 2\ P\ ×\ Lp\ E\ 2\ P)\ :}
\leanline{4}{eLpNorm\ (fun\ ω\ ↦\ weightedPhaseTransform\ a\ b\ (lpPhaseCoe\ P\ S\ ω))\ 2\ P\ =}
\leanline{6}{ENNReal.ofReal\ (weightedPhaseNorm\ a\ b\ S)}
\end{leancode}
\leanlabel{How the Lean proof follows the argument}
\texttt{lp\allowbreak{}Weighted\allowbreak{}Transform} names the mathematical object defined above. Its defining expression appears in the adjacent statement; subsequent identities may unfold this definition directly.




\texttt{lp\allowbreak{}Weighted\allowbreak{}Transform\_\allowbreak{}coe} gives the formal counterpart of this step. Its remaining equalities are discharged by the simplifications displayed in the proof.




\texttt{lp\allowbreak{}Weighted\allowbreak{}Transform\_\allowbreak{}norm} gives the formal counterpart of this step. The main intermediate facts are \texttt{h1}, \texttt{h2} in that order. Named proof applications include \texttt{With\allowbreak{}Lp.\allowbreak{}prod\_\allowbreak{}norm\_\allowbreak{}sq\_\allowbreak{}eq\_\allowbreak{}of\_\allowbreak{}L2}, \texttt{h1.\allowbreak{}trans}; the adjacent proof shows their hypotheses and order of use.




\texttt{e\allowbreak{}Lp\allowbreak{}Norm\_\allowbreak{}weighted\allowbreak{}Phase\allowbreak{}Coe} gives the formal counterpart of this step.
\end{leanbox}

\section{\texorpdfstring{Keep differences and the entire OBABO step tied to their representatives}{Keep differences and the entire OBABO step tied to their representatives}}


The first lemma records the almost-everywhere representative of a phase difference, using the two component subtraction identities. The second expands OBABO into first thermostat, Verlet, and final thermostat, then composes the already proved representative identities for all three. Consequently the lifted nonlinear force update represents exactly the original finite-dimensional \(O\)-\(B\)-A-\(B\)-\(O\) sampler with the same two supplied noise variables. These equalities preserve the joint coupling; a marginal-law equality is never substituted for the common-noise realization.



The representative identities hold for differences and for each affine stage of the OBABO update. Thus the weighted local-error norm is the root-mean-square distance between the coupled vector outputs. This connects the analytic bound with a transport coupling.

\begin{leanbox}{Lean: \texttt{lpPhaseCoe\_\allowbreak{}sub}, \texttt{lp\_\allowbreak{}obaboStep\_\allowbreak{}coe}}
\leansource{\texttt{lpPhaseCoe\_\allowbreak{}sub} \quad KineticLangevin/\allowbreak{}LpWeightedGeometry.lean, lines 85–88}
\begin{leancode}
\leanline{0}{theorem\ lpPhaseCoe\_sub\ (S\ T\ :\ Lp\ E\ 2\ P\ ×\ Lp\ E\ 2\ P)\ :}
\leanline{4}{lpPhaseCoe\ P\ (S\ -\ T)\ =ᵐ[P]\ fun\ ω\ ↦\ lpPhaseCoe\ P\ S\ ω\ -\ lpPhaseCoe\ P\ T\ ω}
\end{leancode}
\leansource{\texttt{lp\_\allowbreak{}obaboStep\_\allowbreak{}coe} \quad KineticLangevin/\allowbreak{}LpWeightedGeometry.lean, lines 92–103}
\begin{leancode}
\leanline{0}{theorem\ lp\_obaboStep\_coe\ (g\ :\ E\ →\ E)\ (M\ :\ ℝ≥0)\ (hg\ :\ LipschitzWith\ M\ g)}
\leanline{4}{(h\ δ\ σ\ :\ ℝ)\ (S\ :\ Lp\ E\ 2\ P\ ×\ Lp\ E\ 2\ P)\ (Ξ\ Ξ'\ :\ Lp\ E\ 2\ P)\ :}
\leanline{4}{lpPhaseCoe\ P\ (obaboStep\ (ULA.lpDrift\ P\ g\ hg)\ h\ δ\ σ\ S\ Ξ\ Ξ')\ =ᵐ[P]}
\leanline{6}{fun\ ω\ ↦\ obaboStep\ g\ h\ δ\ σ\ (lpPhaseCoe\ P\ S\ ω)\ (Ξ\ ω)\ (Ξ'\ ω)}
\end{leancode}
\leanlabel{How the Lean proof follows the argument}
\texttt{lp\allowbreak{}Phase\allowbreak{}Coe\_\allowbreak{}sub} gives the formal counterpart of this step. Named proof applications include \texttt{Prod.\allowbreak{}ext}; the adjacent proof shows their hypotheses and order of use.




\texttt{lp\_\allowbreak{}obabo\allowbreak{}Step\_\allowbreak{}coe} gives the formal counterpart of this step. The main intermediate facts are \texttt{G} in that order.
\end{leanbox}


Earlier results used here: \cite[Section~3.13.2]{LangevinLean2026ULASC}.

\begin{thebibliography}{ULA-LSI}
\bibitem[ULA-SC]{LangevinLean2026ULASC}
LangevinLean Project. \emph{Unadjusted Langevin under Strong Convexity}. LL-01, 2026. Draft report, content revision d34ead06987c. \url{https://odakyildiz.github.io/LangevinLean/docs/pdf/UnadjustedLangevin/FullProof.pdf}.
\end{thebibliography}

\end{document}
